\chapter{Evolutionary Dynamics: A Case Study}
\label{ch:evolution}

The ecosystem of Chapter~\ref{ch:ecosystem} was built to \emph{display}
predator/prey dynamics.  This chapter puts the same kind of machinery to
work \emph{answering a scientific question} --- and in doing so
introduces two ingredients the earlier model lacked: creatures that
\emph{move}, and creatures that carry competing \emph{strategies}
selected over evolutionary time.  The material is adapted from the
author's later study \emph{Veridical Perception and Evolutionary
Dynamics}; here we develop in full the compact mean-field model at its
core, whose every claim can be checked in a few lines of Python, and
describe the spatial simulation in prose, since it is a direct
descendant of \texttt{cadyn.ecosystem}.

\section{The Question: Does Evolution Favor Accurate Perception?}

It is tempting --- and widely assumed among cognitive scientists and AI
researchers --- that natural selection favors \emph{veridical}
perception, perception that reports the world as it actually is.
Accurately registering reality surely helps one survive.  Donald Hoffman
and colleagues challenged this in a 2010 paper, \emph{Natural Selection
and Veridical Perceptions}, arguing that selection optimizes
\emph{fitness}, not truth, and that the two can come apart: a cheap,
fast, systematically \emph{inaccurate} perceptual strategy can outcompete
an accurate but costly one.

The intuition is a matter of timing.  Accurate perception takes
processing --- more neural layers, more computation, more time --- and
neural signals are slow.  In a life-or-death moment, a bird that takes
an extra beat to resolve exactly what it is seeing may already be in the
cat's jaws.  The question is whether this cost is enough to select
\emph{against} veridical perception, and if so, when.

\expand{Hoffman has since developed this line into the \emph{Interface
Theory of Perception}: the thesis that our perceptions are like a
desktop interface --- useful, fitness-tuned icons that hide rather than
reveal the underlying reality, much as a file icon tells you nothing
about voltages on a disk.  His 2019 book \emph{The Case Against Reality}
presents the ``Fitness-Beats-Truth'' theorem, which within a class of
evolutionary-game models shows truth-tuned strategies are driven to
extinction by fitness-tuned ones of equal complexity with probability
approaching one as the world's complexity grows.  The CA experiments of
this chapter are a small, concrete, spatially explicit instance of the
same phenomenon.}

\section{A Mean-Field Model of Strategy Competition}
\label{sec:meanfield}

Consider one species with several competing perceptual strategies.
Strategy~$i$ has a fixed \emph{probability of success} $p_i$ ---
loosely, the chance it survives the encounters of one generation --- and
we take this $p_i$ to be its fitness.  For the running example we use
three strategies with
\[
  p = (p_R,\, p_P,\, p_V) = (0.6,\ 0.75,\ 0.9),
\]
standing for a \textbf{R}andom (non-perceiving), a \textbf{P}robabilistic
(sometimes-perceiving), and a \textbf{V}eridical (always-perceiving)
strategy.  Under the natural reading in which more accurate perception
yields more successful encounters, $V$ has the highest bare fitness.

Now place the strategies on the 2-dimensional toroidal lattice of the
previous chapters, one strategy per cell, with the eight-cell Moore
neighborhood.  The update is deliberately simple and partly stochastic:
at each tick a cell keeps its own strategy with probability equal to its
success rate, and otherwise switches to the majority strategy among its
neighbors.  A high-fitness strategy is ``sticky'' --- its holders rarely
defect --- while low-fitness cells readily convert to whatever surrounds
them.

To analyze this, make one very reasonable mean-field assumption: that the
probability a given neighbor is of type~$i$ equals that type's current
population proportion $s_i$.  Writing $s(t) = (s_1, s_2, s_3)^{\mathsf T}$
for the proportion vector, the dynamics become a matrix iteration
\begin{equation}
  s(t+1) = W(t)\, s(t),
  \label{eq:meanfield}
\end{equation}
where the transition matrix is itself a function of the current state:
\begin{equation}
  W(t) =
  \begin{bmatrix}
    p_1 + (1-p_1)s_1 & (1-p_2)s_1 & (1-p_3)s_1\\
    (1-p_1)s_2 & p_2 + (1-p_2)s_2 & (1-p_3)s_2\\
    (1-p_1)s_3 & (1-p_2)s_3 & p_3 + (1-p_3)s_3
  \end{bmatrix}.
\end{equation}
Each column $j$ sums to 1: a cell of type $j$ either keeps its type
(diagonal term $p_j$) or, with probability $1-p_j$, converts to
type~$i$ with probability $s_i$.  This is like a Markov chain, except
that the transition matrix drifts with the population it governs.

Iterating \eqref{eq:meanfield} from any interior start drives the system
to a vertex of the simplex: the fittest strategy takes over completely
and the rest go extinct.  Figure~\ref{fig:mf-runs} shows this for the
equal start $(\tfrac13,\tfrac13,\tfrac13)$ and the uneven start
$(0.5, 0.3, 0.2)$; in both, $V$ wins.  Notably, $V$ wins the second race
\emph{even though it starts as the smallest population} --- higher
fitness, not a head start, decides the outcome.

\begin{figure}[htbp]\centering
\includegraphics[width=\textwidth]{ch7_meanfield_runs}
\caption{The mean-field strategy dynamics from two initial conditions.
The veridical strategy $V$ ($p=0.9$) drives the random and
probabilistic strategies to extinction, even when (right) it begins as
the least common.}
\label{fig:mf-runs}
\end{figure}

The three-strategy state lives on the 2-simplex --- a triangle whose
corners are the pure strategies and whose interior is every mixture.
Figure~\ref{fig:mf-simplex} plots trajectories from many random starts;
every one flows to the $V$ corner.

\begin{figure}[htbp]\centering
\includegraphics[width=0.7\textwidth]{ch7_simplex}
\caption{Trajectories of the proportion vector on the strategy simplex.
Circles mark starts, stars mark the shared destination: the
veridical vertex.}
\label{fig:mf-simplex}
\end{figure}

\subsection*{A Lemma, and a coordinate-free form}

The matrix $W(t)$ is unwieldy, but the map it defines has a clean
closed form.  Let ``$\circ$'' denote the elementwise (Hadamard) product,
$\vec 1$ the all-ones vector, and $I$ the identity.

\begin{lemma}
The update $s(t+1) = W(t)\,s(t)$ can be written, in any dimension, as
\begin{equation}
  s(t+1) = \Bigl(\, p \circ I \;+\; \vec 1\,\bigl((\vec 1 - p)\circ s\bigr)^{\mathsf T}\Bigr)\circ s,
  \label{eq:lemma}
\end{equation}
equivalently and more concretely as
\begin{equation}
  s(t+1) \;=\; p \circ s \;+\; s\,\bigl\langle \vec 1 - p,\ s\bigr\rangle,
  \label{eq:lemma-scalar}
\end{equation}
where $\langle\,\cdot\,,\cdot\,\rangle$ is the ordinary dot product.
\end{lemma}

\begin{proof}
Suppressing the $t$-dependence and writing $q_j = 1-p_j$, the
$i$\textsuperscript{th} component of $W s$ is
\[
  (Ws)_i = \Bigl(p_i + q_i s_i\Bigr)s_i
           + \sum_{j\neq i} q_j s_i s_j
         = p_i s_i + s_i \sum_{j} q_j s_j
         = p_i s_i + s_i\,\langle \vec 1 - p,\ s\rangle,
\]
which is exactly \eqref{eq:lemma-scalar}; expression \eqref{eq:lemma}
is the same computation arranged as a Hadamard product of the bracketed
matrix (diagonal $p_i$, every row equal to $((\vec 1 - p)\circ s)^{\mathsf T}$)
with the column $s$, followed by the row sum.  Nothing in the derivation
used $n=3$, so the form is valid in every dimension.
\end{proof}

The scalar form \eqref{eq:lemma-scalar} is illuminating: each generation
scales strategy~$i$ by its own fitness $p_i$ and then adds back a share
$s_i$ of the total ``defection mass'' $\langle \vec1 - p, s\rangle$.
Because the surviving-and-switching bookkeeping is exact, the true map
conserves $\sum_i s_i$; in floating-point iteration one renormalizes each
step, which is what ``population \emph{proportions}'' means anyway.  The
module \texttt{cadyn.evolution} provides all three formulations and a
routine \texttt{verify\_lemma} confirming they agree --- for the worked
example and for hundreds of random cases in dimensions 2 through 8.

\subsection*{An empirical theorem}

Simulation makes the following extremely plausible; a formal proof
eluded the original study and we state it as the paper did.

\begin{theorem}[empirical]
Let $s(0)$ be any interior point of the probability simplex
($s_i(0) > 0$ for all $i$), and suppose the fitness vector has a unique
maximum, $p_k > p_j$ for all $j \neq k$.  Then $s(t) \to e_k$ as
$t\to\infty$, where $e_k$ is the standard basis vector with a $1$ in
position $k$.
\end{theorem}

The routine \texttt{cadyn.evolution.dominant\_strategy\_wins} tests one
initial condition; across 500 random starts in three dimensions, and 100
each in dimensions 4, 6, and 10, the unique fittest strategy won every
time --- somewhat stronger evidence than the original paper marshalled.
A natural next step, which the original study flagged as future work, is
to find a replicator system approximating this map and prove convergence
analytically; the replicator equation is the standard bridge from such
population updates to a dynamical system with known convergence theory.

The mean-field model can also \emph{preview} the chapter's central
reversal.  Nothing in it explains why $V$ should be fittest --- that is
an input.  If accurate perception carries a computational cost, we can
model that cost crudely by lowering $V$'s effective success rate.
Figure~\ref{fig:mf-penalty} repeats the equal-start race with $V$'s
fitness cut from $0.9$ to $0.55$, below $P$'s $0.75$: now the
probabilistic strategy sweeps the board and veridical perception is the
one driven extinct.  The mean-field model thus contains both outcomes;
which one obtains depends entirely on whether accurate perception is
worth its price.  The spatial model of the next section supplies that
price from first principles, as a cost in \emph{time}.

\begin{figure}[htbp]\centering
\includegraphics[width=\textwidth]{ch7_penalty}
\caption{Left: with no cost, the veridical strategy wins (as before).
Right: charging veridical perception a computational cost --- modeled
here by lowering its effective fitness below the probabilistic
strategy's --- reverses the outcome, and $V$ goes extinct.}
\label{fig:mf-penalty}
\end{figure}

\section{The Spatial Experiment: Movement, Fitness, and a Time Penalty}
\label{sec:spatial-evolution}

The mean-field model already vindicates Hoffman's logic, but only as
abstractly as his own argument: it \emph{assumes} $V$ has the highest
fitness and then shows the fittest wins.  It says nothing about
\emph{why} veridical perception might fail to be fittest.  For that we
need a model where perception actually costs something, and the natural
vehicle is a spatial ecosystem in the manner of
Chapter~\ref{ch:ecosystem} --- extended, crucially, with
\emph{movement}.

\subsection*{What is added to the Chapter~5 ecosystem}

The second experiment keeps the toroidal lattice, the plant/prey/predator
roles, the stochastic asynchronous updates, and the population and
phase-portrait diagnostics of Chapter~\ref{ch:ecosystem}, and layers on:

\begin{itemize}
\item \textbf{Movement.} Creatures (birds and cats) take random walks: a
  creature moves by copying its state into a vacant neighboring cell and
  vacating its own.  Eating is movement --- a bird eats a plant by
  stepping onto it; a cat eats a bird by pouncing onto its cell.  (The
  author notes that mobile creatures in an evolutionary-game CA appear
  to be novel, tracing the idea to his own 1997 work --- the very book
  this volume revises.)
\item \textbf{Individual fitness, birth, and death.} Each creature
  carries a fitness counter: eating raises it, each idle tick lowers it,
  zero means starvation, and crossing a threshold triggers haploid
  reproduction onto a neighboring empty or plant cell.  Unlike the
  Chapter~5 model, which tracked only population \emph{counts}, here
  fitness is an attribute of each individual.
\item \textbf{Perceptual strategies.} Every bird and every cat carries
  one of three strategies --- $R$ (random walk, no perception), $V$
  (veridical: birds flee perceived cats, cats pounce on perceived
  birds), and $P$ (do $V$ with probability $\rho$, else $R$).  This is
  the mean-field model's $R/P/V$ competition, now embodied in moving,
  eating, breeding animals.
\end{itemize}

Including plants alongside prey and predators yields a
rock-paper-scissors (``Rochambeau'') structure, expressed --- as the
author wryly notes --- through negations: no plants starves the birds,
no birds starves the cats, and no cats lets the birds strip the plants.
That cyclic non-transitivity may be exactly what underlies the dynamic
stability of a broad range of Lotka--Volterra systems; flattening it by
omitting plants obscures the mechanism.

\subsection*{The result, and the reversal}

Run without any timing asymmetry, the spatial model echoes the
mean-field one: the veridical strategies drive the others to extinction,
first the random strategies and then the probabilistic ones, for both
birds and cats.  The surviving bird/cat subsystem settles into the noisy,
damped Lotka--Volterra oscillation familiar from
Chapter~\ref{ch:ecosystem}, spiralling toward an interior Nash
equilibrium; in the reported run the normalized bird/cat proportions
converge near $[0.738,\, 0.262]$ with strikingly small steady-state
variance, and an FFT of the population series reveals a dominant cycle of
about 18 generations.  So far, ``veridical perception wins again.''

The decisive experiment introduces the cost of perception as a cost in
\emph{time}.  In the baseline simulation, cats happen to move first ---
they perceive and act before birds do.  To penalize veridical
perceivers for the processing their accuracy requires, the veridical
creatures are made to move \emph{last}: they pay for their better
picture of the world with a beat of delay, exactly the slow-neuron cost
of the opening question.  Under this single change the balance shifts.
The probabilistic strategy, now no longer strictly dominated, competes
with the veridical one on even terms and is no longer driven to
extinction; a slow-but-accurate strategy can be beaten by a
faster, less accurate one.  A pair of thought experiments makes the
mechanism vivid: a bird that resolves ``cat'' a moment too late is
already caught, and a cat that resolves ``bird'' too late goes hungry
--- and repeated, either lag is lethal to the strategy that carries it.

\expand{This is the CA analogue of Hoffman's claim and, more broadly, a
lesson machine learning has relearned: depth is not free.  Because
neural signals are slow, biological brains run \emph{shallow} networks
compared with the very deep ones routinely trained today; a perceptual
system optimized purely for accuracy, with no regard to latency, is not
the system evolution builds.  ``To have survival value, perceptions must
be fast'' --- a constraint now visible in real-time robotics and
edge-inference, where a slightly-worse-but-faster model routinely beats a
slightly-better-but-slower one.}

\section{An Emergent Curiosity: Fitness Collapse Before Extinction}

One observation from the spatial runs deserves flagging because it was
not sought and is not obviously an artifact.  Tracking the \emph{mean
fitness} of each sub-population over time, the author found that a
strategy's mean fitness tends to \emph{oscillate wildly in the moments
before that strategy goes extinct}, and that the overall mean fitness of
the whole ecosystem \emph{steps down} with each extinction event.  Both
effects invite the question the original study left open: are they
quirks of this particular CA, or do they echo something real ---
critical-slowing-down and variance-inflation of the kind studied as
\emph{early-warning signals} of tipping points in real ecosystems?  The
CA is cheap enough to make this a natural thing to probe systematically,
and the tools of Chapter~\ref{ch:ecosystem} (population and phase-space
diagnostics) apply directly.

\section{Reflection}

Two models, one abstract and one spatially explicit, independently
support Hoffman's conclusion: veridical perception is favored by
selection \emph{until} its computational cost is reckoned, after which a
faster, less accurate strategy can win.  The mean-field model is fully
reproduced and checked in \texttt{cadyn.evolution}; the spatial model,
being a direct extension of \texttt{cadyn.ecosystem} with movement and
per-individual strategies, is left as a described experiment and an
invitation.  It closes a loop with Chapter~\ref{ch:ecosystem}: the
predator/prey CA that began as a way to \emph{watch} population dynamics
turns out, with movement and strategies added, to be a laboratory for one
of the sharper questions in the science of the mind --- whether the world
we perceive is the world as it is, or merely the world as it pays to see.
